\section{Returning to Patrol}\label{sec:patrol}
The same retained entry now separates the theorem's two conditions. We
first show a winning source policy and an immediate safe response. We then
show that \emph{every} safe strategy either fails its recurring obligation
or reaches a state with no first-change response.

\subsection{The positive continuation}
The route in Fig.~\ref{fig:patrol-grid} waits two cells below each checkpoint
until its \emph{next} blocked flag is false, approaches through row three,
visits and retreats to row two. After the three visits it returns via $O$
to $H$, dwells until the counter reaches three, and returns through $(4,1)$,
$(4,2)$ and row two to the first standoff. All moves are unit cardinal steps.

The source rules justify each part of this policy. Positions of checkpoints
are fixed. When the robot enters a checkpoint's one-cell neighborhood its
flag is false; neighborhood stability keeps it false through the visit and
retreat. The rest of the route avoids the checkpoints and obeys Dwell.
ENV fairness ends each wait, so the visits recur.

For charging, the counter increments, capped at three, when the \emph{current}
position is $H$ and is zero on every other successor. The monitor is idle,
pending or failed. An idle monitor becomes pending when charging is unavailable,
the current counter is not three and the robot is off the charger. A pending
monitor first tests whether the counter is three, entering the absorbing failed
state if so; otherwise it clears to idle when the current position is $H$.
It remains pending otherwise. The complete SYS justice list requires infinitely many idle states
and visits to each of the three checkpoints and $O$ (five predicates). The complete source equations
are in Supplements~S and AB.6.

First arrival at $H$ has counter zero. The pending request clears on the
first dwell edge, before the counter reaches three, and cannot retrigger at
$H$. The source requires charging to become available on an edge whose
current counter is three. Departure therefore carries counter three, an idle
monitor and available charging. On the next off-charger edge the current
counter prevents retriggering while it resets. The monitor never fails and
is idle infinitely often. Thus the same policy meets both safety and progress:
$r\in W$.

At $r$, any typed next ENV input permits staying at $O$, keeping counter zero
and the request pending. Dwell is inactive and all current blocked flags are
false. Hence $C(r)$ also holds. What is missing is the ability to maintain
$C$ while winning.

\subsection{Why every safe strategy loses that option}
Let ENV keep the displayed layout, every checkpoint open and charging
unavailable until the first arrival at $H$. This is source-legal; keeping every checkpoint open also satisfies all
three ENV justices. Figure~\ref{fig:patrol-motivation} divides all safe
strategies into two cases.

\emph{Never reach the charger.} Every safe response keeps the counter zero
and the request pending: the updates inspect the current off-charger position.
The resulting ENV-fair play never clears the request, violating SYS justice.
Thus a source-winning strategy cannot choose this branch forever.

\emph{First reach the charger.} The arrival still has counter zero and a
pending request. On the next edge ENV legally closes only distant checkpoint
$A$. Dwell keeps the robot at $H$; the counter becomes one and the request
clears. Call this state $x$. Now a typed first-change input places that
already blocked checkpoint at $H$. Dwell requires $p'=H$ and Avoidance
requires $p'\ne H$. No response exists, so $x\notin C$.
The contradiction occurs \emph{after the request clears}, while the dwell
is still unfinished.

The two branches cover all safe responses, not just the route we displayed.
Consequently $r\in W\cap C$ but $r\notin W_C$. Theorem~\ref{hpb:exact}
excludes $(r,0)$ from the complete literal game, even though the source can
continue from the same entry. The loss is caused by having to keep a response
available until an adversary chooses the first change.

This proves the loss for complete literal relations. Tool exports may
retain auxiliary duties after the helper changes. We extend the argument
before applying it to Patrol's actual export.
\subsection{If safety duties remain after the change}
A syntactic translator may leave counter or monitor constraints unguarded.
Let $U\subseteq D$ be exactly the states admitting a SYS strategy that
satisfies all remaining safety duties along every legal helper-true ENV
prefix. ENV dead ends are winning vacuously. Write $\rho_e^{01}$ and
$\rho_s^{01}$ for the complete encoded ENV and SYS safety relations with
$b=0,b'=1$. The exact safe-entry condition is
\begin{equation}
\begin{aligned}
 F(s)\Longleftrightarrow{}&\forall e'\in E:\ \rho_e^{01}(s,e')\Rightarrow{}\\
 &\exists y'\in Y:\ \rho_s^{01}(s,e',y')\land(e',y')\in U.
\end{aligned}
\label{hpb:F}
\end{equation}
Thus every legal first-change input must admit a safe successor in $U$.
Define $W_F$ by replacing $C$ with $F$ in the maintained-response
definition of $W_C$, including its initial-state requirement.

\begin{proposition}[Preservation with retained duties]\label{hpb:retained}
Suppose the safety relations on $0\to0$ edges agree with the source.
Apart from the added $\G\F\neg b$ ENV assumption, the ENV and SYS justice
lists agree with the source at helper false and may be arbitrary at
helper true. The helper is one-way, its next value is observed before
SYS responds, and ENV legality has no next-SYS dependence.
Then the helper-true strict-winning region is $U$ and the helper-false
region is $W_F$. Whole-region preservation holds exactly when $W\subseteq F$.
\end{proposition}
\begin{proof}
At helper true, persistence makes $\G\F\neg b$ false on every infinite
legal play. Strict winning is therefore exactly safety winning, namely $U$.

Restrict a winning helper-false strategy to source-legal histories where
the helper stays false. Their obligations coincide with the source.
Every such prefix, including the initial state, has a winning residual
with its reached memory: required safety has held and the B\"uchi
objectives are prefix-independent. Every legal first change must
therefore receive a safe response into $U$. This same restricted strategy
wins the source while maintaining $F$, proving membership in $W_F$.

Conversely, follow one source-winning strategy witnessing $W_F$, with
its evolving memory, until a first change. Then $F$ supplies a safe response
into $U$; continue with a safety-winning strategy on $U$. Without a change,
source obligations hold unchanged. An encoded-illegal input releases
strict duties on that edge and thereafter. This proves both region equalities.

Finally, $W_F\subseteq W\cap F$ gives necessity for whole-region
preservation. If $W\subseteq F$, every source-legal prefix of a source-winning
strategy has a winning residual, so every reached state lies in
$W\subseteq F$. The same strategy witnesses $W_F$, proving sufficiency.
\end{proof}

Supplement~L.1 gives an example with two source states, one always-legal input
value, a Boolean SYS output, and true source initial and justice predicates:
source steps and first-change responses must set the next output to true.
Every subsequent helper-true transition requires both current and next
outputs to be false. Thus both source states win and admit an immediate
safe first-change response, but that response enters a state with no safe
next response, so neither state survives the encoding.
Supplement~L.1 gives the complete relational formulation and fixed-point
construction. For complete literal relations, $U=D$ and $F=C$; these
identities must be checked for an actual export.


\subsection{The actual export loses the same entry}
For the 4,800 current states with the displayed layout, a completed native
reduced calculation finds \textbf{3,280 source winners and 1,104 survivors},
agreeing with a separate explicit all-state audit. The fixed layout limits
the counted current states, not the typed next inputs: the first-change
test still permits a checkpoint to move to any typed position. The reduction enforces
maintenance of $F$ while winning (Supplement~L.1). The same retained entry
is excluded by the actual syntactic export, by the following relation checks.

The export retains auxiliary updates, so Proposition~\ref{hpb:retained}
applies. Complete typed checks establish source agreement before the change
and unchanged helper-false ordered justices. At the first change every typed
ENV input is legal and the source SYS duty remains. Afterwards ENV is
unrestricted; the remaining counter/monitor updates permit every typed robot
position with a typed auxiliary update, independently of the next ENV values.
Thus $U=D$ and $F=C$, and the all-response argument excludes $(r,0)$ from
the export's strict-winning region. This is an export-based strict deduction, not an observed public Kind
refusal; the direct full-helper attempts remain UNKNOWN (Supplements~U and AB.4).

\subsection{What the retained-state execution shows}
The recorded update reused the legal policy under the old and then new
ENV model at the same retained query, preserving plant/controller objects,
source fields, monitors and private phase. It completed the pending request
and entered a recurring justice-satisfying cycle. The separate source-kind
existence argument follows in Section~\ref{sec:source-kind}.
